\minitopic{怀疑论与反运气理论共享一种语法}怀疑者说：在某个与实际经验不可区分的情形中，命题为假；主体不能排除该情形；因此主体不知道。反运气理论也考察非实际情形，只是方向相反：在主体以同样方法形成信念的近邻情形中，信念是否容易为假。两类论证都需要回答哪些情形相关、保持哪些条件不变、距离如何比较。差别在于，怀疑论常从一个极端替代情景挑战知识范围，反运气理论则从局部脆弱性诊断一次成功为何不够资格成为知识。

\minitopic{逻辑可能不等于认识上相关}若任何逻辑可能都必须排除，我们在看见桌子时还要排除完美模拟、瞬间复制和不可检测的宇宙改变，求知实践将无法开始。可错论因此允许剩余错误风险，但不能简单说“日常不在意的可能都不相关”。相关性至少可来自三处：当前证据中的异常使某替代情景得到具体支持；环境统计结构使错误在相似场景中频繁发生；任务本身要求排除某类错误，例如法庭鉴定文件真伪。相关性判断必须能够被公开说明，否则它只是在知识直觉之后补写理由。

\minitopic{闭合争论揭示理论成本}保留闭合可以解释知识怎样通过已知蕴涵扩展，也让怀疑论从普通命题推进到反怀疑命题。拒绝闭合可以阻断推进，却会允许主体知道“这是斑马”而不知道“它不是伪装骡子”，并损害证明的认识作用。较温和的做法给闭合加入主体掌握蕴涵、完成推理和适当基于推理形成信念等条件，但怀疑压力仍不会完全消失。因而回应怀疑论不能只挑一条前提删去，还要报告删除后普通知识、推理扩展与反省知识分别受到什么影响。

\minitopic{近邻世界必须固定方法}敏感性问：若$p$为假，主体是否仍会相信$p$；安全性问：主体以当前方式相信$p$时，信念是否容易为假。比较时若同时改变证据、注意、信念方法与环境，就能任意制造想要的结论。彩票持有人基于基础率相信自己会输；在中奖世界中他仍会这样相信，因此不敏感。一个直接看见正常谷仓的人在假立面密布的环境中实际相信真命题，却在近邻同方法情形中很容易出错，因此不安全。方法固定让两个条件真正检验信念与真值的反事实联系，而不是检验故事作者的想象偏好。

\begin{argumentbox}
处理错误可能时依次问四题：该可能是否得到具体证据支持；它在当前环境中是否容易实现；主体的方法能否在其实现时改变信念；任务或制度是否特别要求排除它。四题给出的不是自动算法，而是一份可争论、可修正的相关性说明。
\end{argumentbox}

\minitopic{能力理论补充“为什么成功”}安全性可以说明信念不容易为假，却未必说明成功属于谁。高度可靠的自动设备替一个完全不理解输出的人给出答案，可能满足某些反事实条件，却不一定构成他的能力成就。德性理论要求成功在适当程度上归功于认知能力\parencite{sosa2007,zagzebski1996}。但现代知识往往由个人、仪器、软件和组织程序共同产生，因此归功不能预设单一个体。好的能力分析要指出每一层负责什么、错误如何被发现，以及若移除某项能力，成功是否仍稳定出现。

\minitopic{怀疑可以被制度化而不必无限化}科学复现、双人复核、来源签名和异议程序都把特定怀疑转化为可执行检查。它们没有证明一切错误不可能，而是在高风险节点选择最相关失败模式，分配核查责任，并保存修正记录。这说明怀疑的实践价值不在持续否定，而在改善认识环境。一个理论若只能告诉主体“仍有可能错”，却不能区分应继续调查与可以暂时行动的情形，就尚未完成认识论任务。

\minitopic{阶段性结论}知识不要求排除全部可能，却要求主体与环境之间存在足以抵抗相关错误的联系。怀疑论帮助我们看见未经辩护的停止点，反运气理论帮助我们检验成功的脆弱性，能力理论再追问成功归属于谁。把三者连接起来，知识边界便不再由“绝对确定”或“看起来正常”两个极端决定，而由可说明的相关性、反事实稳定性和能力贡献共同约束。
